\documentclass[12pt]{amsart}
\usepackage{amsmath,amssymb,amsthm}
\usepackage[margin=1in]{geometry}
\usepackage{hyperref}

%% ── Theorem environments ───────────────────────────────────────────────────
\newtheorem{theorem}{Theorem}[section]
\newtheorem{lemma}[theorem]{Lemma}
\newtheorem{proposition}[theorem]{Proposition}
\newtheorem{definition}[theorem]{Definition}
\theoremstyle{remark}
\newtheorem{remark}[theorem]{Remark}
\newtheorem{openproblem}{Open Problem}

%% ── Macros ─────────────────────────────────────────────────────────────────
\newcommand{\R}{\mathbb{R}}
\newcommand{\Tcal}{\mathcal{T}}
\newcommand{\lH}{\lambda_H}
\newcommand{\Qsix}{Q_6}
\newcommand{\jstar}{j_*(t)}
\newcommand{\ip}[2]{\langle #1,\, #2 \rangle}
\newcommand{\norm}[2]{\|#1\|_{#2}}
\newcommand{\eps}{\varepsilon}

\title{Global Regularity of the Three-Dimensional Incompressible\\
Navier--Stokes Equations via Self-Adaptive Spectral Damping}
\author{Jonathan Simons}
\address{Prime Field Technologies LLC}
\email{jonathansimons357@proton.me}
\date{2026}

\begin{document}
\maketitle

%% ─────────────────────────────────────────────────────────────────────────
\begin{abstract}
We construct a self-adaptive augmentation of the three-dimensional
incompressible Navier--Stokes equations on $\R^3$ enforcing global
regularity through two mechanisms: (i)~an adaptive coherence viscosity
$\lH[u]$ defined by an auxiliary elliptic problem, and (ii)~a
spectrally-localized damping operator $\Qsix[u]$ supported on the
dynamically chosen dominant Fourier shell
$j_*(t)=\operatorname{argmax}_j\{2^{2j}\|\Delta_j u\|_{L^2}^2\}$.
The dominant-shell construction makes the coercive estimate immediate
and closes a gap present in average-frequency formulations.

We prove global well-posedness in $H^m$ ($m\ge 3$) via Galerkin methods
invoking Carath\'eodory's existence theorem for the finite-dimensional
projections, uniform energy bounds, vorticity control, and attractor
convergence. All constants are explicit and tracked from Sobolev and
Ladyzhenskaya inequalities in $\R^3$.

The central structural result is the \emph{Spectral Dichotomy}
(Lemma~\ref{lem:dichotomy}). In Case~I (concentration) the energy
inequality factors cleanly via the dominant-shell bound. In Case~II
(dispersal) a complete Bony paraproduct argument with explicit
Coifman--Meyer remainder estimate recovers the standard cubic bound.
The $H^m$ coercivity of $\Qsix$ follows from the commutativity of
$D^\alpha$ and Fourier multipliers, eliminating uncontrolled
high-frequency coefficients from the energy estimate.

The de-augmentation problem reduces precisely to the
\textbf{Spectral Non-Dispersal Condition~[SND]:} under~[SND]
the augmented solutions converge strongly to a globally smooth solution
of the classical equations via Aubin--Lions compactness. [SND] is
physically natural, mathematically precise, and currently open.
A proof of~[SND] combined with Theorem~\ref{thm:A} constitutes a
complete proof of global regularity for classical 3D Navier--Stokes.
\end{abstract}

\tableofcontents

%% ─────────────────────────────────────────────────────────────────────────
\section{Introduction}
%% ─────────────────────────────────────────────────────────────────────────

The global regularity of the three-dimensional incompressible
Navier--Stokes equations is a Clay Mathematics Institute Millennium
Prize Problem. Leray~\cite{Ler34} established global weak solutions;
whether they remain smooth for all time is unresolved.  The fundamental
obstruction is supercriticality: vortex stretching $(\omega\cdot\nabla)u$
can in principle drive enstrophy to infinity, and the
Beale--Kato--Majda criterion~\cite{BKM84} identifies
$\int_0^T\|\nabla u\|_{L^\infty}\,dt$ as the critical blow-up quantity.

This paper constructs an augmented system with two explicit,
self-adaptive components.
\begin{enumerate}
\item \emph{Adaptive coherence viscosity} $\lH[u]$
  (Section~\ref{sec:lH}): a scalar, time-dependent additional viscosity
  whose coefficient is determined by an auxiliary elliptic problem,
  increasing dissipation precisely where gradients concentrate.
\item \emph{Dominant-shell operator} $\Qsix[u]$
  (Section~\ref{sec:Q6}): a Fourier multiplier supported on the three
  dyadic shells surrounding the dominant gradient shell $\jstar$.
  This provides sign-definite dissipation at the most energetically
  active frequency scales.
\end{enumerate}

The key structural advance over prior formulations is replacing the
average-frequency scale
$\Lambda(t)=\|\nabla u\|_{L^2}/\|u\|_{L^2}$---which does not guarantee
shell concentration---with the dominant-shell index
$j_*(t)=\operatorname{argmax}_j\{2^{2j}\|\Delta_j u\|_{L^2}^2\}$,
making the coercive estimate an immediate consequence of the
construction.

%% ─────────────────────────────────────────────────────────────────────────
\section{Classical Framework and Notation}
\label{sec:classical}
%% ─────────────────────────────────────────────────────────────────────────

We work on $\R^3$.  For $s\ge 0$, $H^s(\R^3)$ carries the norm
$\|u\|_{H^s}^2=\int(1+|\xi|^2)^s|\hat u(\xi)|^2\,d\xi$.
The Leray projector $P$ onto divergence-free fields satisfies
$\widehat{Pv}(\xi)=v(\xi)-(\xi\cdot v(\xi)/|\xi|^2)\xi$.
The Littlewood--Paley projector $\Delta_j$ acts on the dyadic shell
$S_j=\{\xi:2^{j-1}\le|\xi|<2^j\}$; set $S_j(u)=\sum_{k\le j}\Delta_ku$.

\medskip\noindent\textbf{Key quantities.}
\begin{align}
  j_*(t) &:= \operatorname{argmin}\!\Bigl\{j\in\mathbb{Z}:
    2^{2j}\|\Delta_j u(t)\|_{L^2}^2 =
    \max_{k\in\mathbb{Z}} 2^{2k}\|\Delta_k u(t)\|_{L^2}^2\Bigr\},
    \label{eq:jstar}\\
  \Tcal(u) &:= \sum_{|j-j_*(t)|\le 1}
    2^{2j}\|\Delta_j u(t)\|_{L^2}^2,
    \label{eq:T}\\
  E(t) &:= \tfrac{1}{2}\|u(t)\|_{H^1}^2.
    \label{eq:E}
\end{align}

The classical system is
\begin{equation}\label{eq:NS}
  \partial_t u + (u\cdot\nabla)u + \nabla p = \nu\Delta u,
  \quad \nabla\cdot u = 0,
  \quad u(\cdot,0)=u_0\in H^1(\R^3).
\end{equation}
The classical energy identity
$\frac{d}{dt}E_L(t)=-\nu\|\nabla u\|_{L^2}^2$
ensures global dissipation but does not control gradient concentration.
The BKM criterion~\cite{BKM84}: if
$\int_0^T\|\nabla u\|_{L^\infty}\,dt<\infty$ then
$u\in C^\infty(\R^3\times[0,T])$.

%% ─────────────────────────────────────────────────────────────────────────
\section{Adaptive Coherence Viscosity}
\label{sec:lH}
%% ─────────────────────────────────────────────────────────────────────────

\subsection{Auxiliary elliptic problem}
Define $\psi[u]$ as the unique solution (decaying at infinity) of
\begin{equation}\label{eq:psi}
  -\Delta\psi = \nabla\cdot(|u|^2u)\quad\text{on }\R^3.
\end{equation}
By Lax--Milgram this is well-posed for $u\in H^2(\R^3)$, with
$\|\nabla\psi\|_{L^2}\le C_S\||u|^2u\|_{L^2}$.
For $u\in H^m$, $m\ge 2$, Sobolev embedding gives
$\|\nabla\psi\|_{L^2}\le\|u\|_{L^\infty}^2\|u\|_{L^2}\le C\|u\|_{H^m}^3$.

\subsection{Definition}
For $\alpha>0$ set
\begin{equation}\label{eq:lH}
  \lH[u]:=\alpha\|\nabla\psi[u]\|_{L^2}^2.
\end{equation}
Properties: (a)~$\lH\ge 0$; (b)~$\nu+\lH\ge\nu>0$ uniformly;
(c)~$\lH\le C\alpha\|u\|_{H^m}^6$.
Because $\lH$ is spatially scalar, no commutator terms arise in $H^m$
energy estimates.

%% ─────────────────────────────────────────────────────────────────────────
\section{The $Q_6$ Operator}
\label{sec:Q6}
%% ─────────────────────────────────────────────────────────────────────────

\subsection{Dominant-shell definition}

\begin{remark}
The average-frequency scale $\Lambda(t)=\|\nabla u\|_{L^2}/\|u\|_{L^2}$
is the $L^2$-weighted mean of $|\xi|$.  A spectral measure can have its
mean at $\Lambda$ while being split between $|\xi|\ll\Lambda$ and
$|\xi|\gg\Lambda$, so no fixed fraction of gradient energy need sit in
any annulus near $\Lambda$.  The dominant-shell construction avoids this
entirely by targeting the shell of maximum gradient energy directly.
\end{remark}

\begin{remark}[Measurability of $j_*(t)$]
The index $j_*(t)$ is integer-valued and in general discontinuous in
$t$---it jumps when two dyadic shells swap dominance.  It is however
measurable: on any interval where $u$ is continuous in $H^1$, each
function $t\mapsto 2^{2j}\|\Delta_j u(t)\|_{L^2}^2$ is continuous for
fixed $j$, so $j_*(t)$ is piecewise constant and hence measurable.
The Galerkin well-posedness argument in Theorem~\ref{thm:gwp} Step~1
invokes Carath\'eodory's existence and uniqueness theorem~\cite{Hartman}:
the Galerkin ODE has right-hand side $f(t,u_n)$ that is (i)~locally
Lipschitz in $u_n$ for each fixed $t$, and (ii)~measurable and locally
bounded in $t$ for fixed $u_n$.  These are precisely the Carath\'eodory
conditions.
\end{remark}

Let $\varphi\in C^\infty_c\bigl((1/2,2)\bigr)$ with $\varphi=1$ on
$[3/4,3/2]$, $0\le\varphi\le 1$.  Define the Fourier multiplier symbol
\begin{equation}\label{eq:symbol}
  m_{j_*}(\xi,t):=|\xi|^2\,\varphi\!\left(2^{-j_*(t)}|\xi|\right),
\end{equation}
and the operator
\begin{equation}\label{eq:Q6}
  \Qsix[u]:=-\gamma\,P\!\left[m_{j_*}(D,t)\,u\right],\quad\gamma>0.
\end{equation}
The support of $m_{j_*}(\cdot,t)$ is contained in the active band
$B(t)=S_{j_*-1}\cup S_{j_*}\cup S_{j_*+1}$.

\subsection{Lemma 1: Coercive dissipation and $H^m$ control}

\begin{lemma}[Properties of $\Qsix$]\label{lem:Q6}
Let $u\in H^m(\R^3)$, $m\ge 3$, divergence-free.  Then the following
hold for a.e.\ $t$\textup{:}
\begin{enumerate}
  \item[\textup{(i)}]
    \textbf{Divergence-free:} $\nabla\cdot\Qsix[u]=0$.
  \item[\textup{(ii)}]
    \textbf{Sign-definite:} $\ip{\Qsix[u]}{u}_{L^2}\le 0$.
  \item[\textup{(iii)}]
    \textbf{Coercive lower bound:}
    $\ip{\Qsix[u]}{u}_{L^2}\le -\tfrac{\gamma}{4}\Tcal(u)$.
  \item[\textup{(iv)}]
    \textbf{$H^m$ coercivity:} For all $|\alpha|\le m$,
    \[
      \ip{D^\alpha\Qsix[u]}{D^\alpha u}_{L^2}\le 0.
    \]
  \item[\textup{(v)}]
    \textbf{$H^{m-1}$ bound:}
    $\|\Qsix[u]\|_{H^{m-1}}\le 4\gamma\cdot 2^{2j_*}\|u\|_{H^{m-1}}$.
\end{enumerate}
\end{lemma}

\begin{proof}
(i)~$\nabla\cdot P[v]=0$ for all $v$.

(ii)--(iii)~By Plancherel and $P$ self-adjoint (with $Pu=u$ since $u$
is divergence-free):
\[
  \ip{\Qsix[u]}{u}_{L^2}
  =-\gamma\int_{\R^3}m_{j_*}(\xi,t)|\hat u(\xi)|^2\,d\xi\le 0,
\]
since $m_{j_*}=|\xi|^2\varphi\ge 0$ pointwise.  On $S_{j_*}$,
$\varphi=1$ and $|\xi|^2\ge 2^{2j_*}/4$.  On adjacent shells
$|j-j_*|=1$, $\varphi\ge 3/4$.  Summing over the active band:
\[
  \ip{\Qsix[u]}{u}_{L^2}
  \le-\frac{\gamma}{4}\sum_{|j-j_*|\le 1}2^{2j}\|\Delta_ju\|_{L^2}^2
  =-\frac{\gamma}{4}\Tcal(u).
\]

(iv)~Since $D^\alpha$ and $m_{j_*}(D,t)$ are both Fourier multipliers,
they \emph{commute exactly}: $D^\alpha m_{j_*}(D,t)=m_{j_*}(D,t)D^\alpha$.
Therefore
\[
  \ip{D^\alpha\Qsix[u]}{D^\alpha u}_{L^2}
  =-\gamma\int m_{j_*}(\xi,t)|\widehat{D^\alpha u}(\xi)|^2\,d\xi\le 0.
\]
No high-frequency coefficient $2^{4j_*}$ appears; the commutativity
is exact.

(v)~$|m_{j_*}|\le 4\cdot 2^{2j_*}$ on its support.
\end{proof}

\begin{remark}
Part~(iv) is the key structural improvement over prior formulations.
Because $\Qsix$ is a Fourier multiplier, testing the $H^m$ equation
against $D^\alpha u$ moves the $\Qsix$ contribution entirely to the
left side as a non-negative dissipative term.  No $2^{4j_*}$ growth
coefficient appears on the right.
\end{remark}

%% ─────────────────────────────────────────────────────────────────────────
\section{The Spectral Dichotomy}
\label{sec:dichotomy}
%% ─────────────────────────────────────────────────────────────────────────

\begin{lemma}[Spectral Dichotomy]\label{lem:dichotomy}
Let $u\in H^m(\R^3)$, $m\ge 3$, be divergence-free.
Fix $\sigma\in(0,1)$.  Exactly one of the following holds.

\medskip\noindent
\textbf{Case~I (Concentration):} $\Tcal(u)\ge\sigma\|\nabla u\|_{L^2}^2$.
In this case
\begin{equation}\label{eq:caseI}
  \frac{d}{dt}E(t)+\frac{\nu}{2}\|\Delta u\|_{L^2}^2
  \le\Tcal(u)\!\left[C_4 E(t)-\frac{\gamma}{4}\right],
\end{equation}
where $C_4>0$ depends on $\nu$, the Sobolev constants in $\R^3$, and
the initial energy $E(0)$.

\medskip\noindent
\textbf{Case~II (Dispersal):} $\Tcal(u)<\sigma\|\nabla u\|_{L^2}^2$.
In this case
\begin{equation}\label{eq:caseII}
  \frac{d}{dt}E(t)+\frac{\nu}{2}\|\Delta u\|_{L^2}^2
  \le C_4 E(t)^3.
\end{equation}
\end{lemma}

\begin{proof}
\textbf{Preliminary (standard 3D nonlinear bound).}
For divergence-free $u\in H^1(\R^3)$, Sobolev embedding
$H^1\hookrightarrow L^6$ (constant $C_S$) and Ladyzhenskaya's
interpolation in $\R^3$,
\[
  \|\nabla u\|_{L^3}\le C\|\nabla u\|_{L^2}^{1/2}\|\Delta u\|_{L^2}^{1/2},
\]
give
\[
  |\ip{(u\cdot\nabla)u}{-\Delta u}|
  \le\|u\|_{L^6}\|\nabla u\|_{L^3}\|\Delta u\|_{L^2}
  \le C_SC\|\nabla u\|_{L^2}^{3/2}\|\Delta u\|_{L^2}^{3/2}.
\]
Young's inequality ($ab\le\frac{\nu}{2}a^{4/3}+C_\eps b^4$ with
$a=\|\Delta u\|_{L^2}^{3/2}$, $b=\|\nabla u\|_{L^2}^{3/2}$,
$\eps=\nu/2$) yields
\begin{equation}\label{eq:prelim}
  |\ip{(u\cdot\nabla)u}{-\Delta u}|
  \le\frac{\nu}{2}\|\Delta u\|_{L^2}^2
  +C_3\frac{\|\nabla u\|_{L^2}^6}{\nu^3},
\end{equation}
with explicit constant $C_3=(C_SC)^4/(2\nu^3)$.
Since $\|\nabla u\|_{L^2}^2\le 2E(t)$,
\begin{equation}\label{eq:C4raw}
  C_3\frac{\|\nabla u\|_{L^2}^6}{\nu^3}\le\frac{8C_3}{\nu^3}E(t)^3.
\end{equation}

\medskip\noindent\textbf{Case~I.}
Since $\Tcal(u)\ge\sigma\|\nabla u\|_{L^2}^2$, we have
$\|\nabla u\|_{L^2}^2\le\Tcal(u)/\sigma$.
Apply this directly to the preliminary bound:
\[
  C_3\frac{\|\nabla u\|_{L^2}^6}{\nu^3}
  = \frac{C_3}{\nu^3}\|\nabla u\|_{L^2}^2\cdot\|\nabla u\|_{L^2}^4
  \le \frac{C_3}{\nu^3}\cdot\frac{\Tcal(u)}{\sigma}
     \cdot\|\nabla u\|_{L^2}^4.
\]
Under the bootstrap $E(t)\le M:=2E(0)$:
$\|\nabla u\|_{L^2}^2\le 2E(t)$ and $\|\nabla u\|_{L^2}^2\le 2M$,
so $\|\nabla u\|_{L^2}^4\le 2E(t)\cdot 2M=4ME(t)$.
Therefore
\begin{equation}\label{eq:caseI_key}
  C_3\frac{\|\nabla u\|_{L^2}^6}{\nu^3}
  \le\frac{4C_3 M}{\sigma\nu^3}\,\Tcal(u)\cdot E(t).
\end{equation}
Setting $C_4:=4C_3M/(\sigma\nu^3)$ (which depends on $\nu$, $C_S$,
$\sigma$, and $E(0)$ through $M=2E(0)$), inequality~\eqref{eq:caseI_key}
gives
\[
  C_4 E^3 - \frac{\gamma}{4}\Tcal(u)
  \le C_4\Tcal(u)\cdot E - \frac{\gamma}{4}\Tcal(u)
  = \Tcal(u)\!\left[C_4 E - \frac{\gamma}{4}\right].
\]
This is the factored form~\eqref{eq:caseI}.

\medskip\noindent\textbf{Case~II (Bony paraproduct expansion).}
When $\Tcal(u)<\sigma\|\nabla u\|_{L^2}^2$, the spectral mass is
dispersed and $\Tcal(u)$ cannot dominate.  Decompose via the Bony
paraproduct~\cite{Bon81}:
\[
  (u\cdot\nabla)u = T_u(\nabla u)+T_{\nabla u}u+R(u,\nabla u),
\]
where $T_f g:=\sum_j S_{j-2}(f)\Delta_j g$ and
$R(f,g):=\sum_j\sum_{|k-j|\le 1}\Delta_j f\cdot\Delta_k g$.

\smallskip\noindent\textit{Low-high term $T_u(\nabla u)$.}
Testing against $-\Delta u$:
\[
  |\ip{T_u(\nabla u)}{-\Delta u}|
  \le\sum_j\|S_{j-2}u\|_{L^\infty}
     \|\nabla\Delta_j u\|_{L^2}\|\Delta\Delta_j u\|_{L^2}.
\]
Since $\|S_{j-2}u\|_{L^\infty}\le\|u\|_{L^\infty}\le C\|u\|_{H^m}$
(Sobolev, $m\ge 3>3/2$), and by Bernstein
$\|\nabla\Delta_j u\|_{L^2}\le 2^j\|\Delta_j u\|_{L^2}$,
summing over $j$ with Cauchy--Schwarz gives
\[
  |\ip{T_u(\nabla u)}{-\Delta u}|
  \le C\|u\|_{H^m}\|\nabla u\|_{L^2}\|\Delta u\|_{L^2}
  \le\frac{\nu}{8}\|\Delta u\|_{L^2}^2 + C(\nu,\|u\|_{H^m})\|\nabla u\|_{L^2}^4.
\]
Under the $H^m$ bootstrap this is absorbed into $C_4E^3$.

\smallskip\noindent\textit{High-low term $T_{\nabla u}u$.}
By the same argument with low-frequency $\nabla u$ and high-frequency
$u$ interchanged:
$|\ip{T_{\nabla u}u}{-\Delta u}|\le\tfrac{\nu}{8}\|\Delta u\|_{L^2}^2+C_4E^3$.

\smallskip\noindent\textit{High-high remainder $R(u,\nabla u)$.}
The Coifman--Meyer theorem~\cite{CM78} applied to the bilinear operator
with symbol smooth away from the origin gives
\[
  \|R(u,\nabla u)\|_{L^2}\le C\|u\|_{L^6}\|\nabla u\|_{L^3}
  \le CC_S\|\nabla u\|_{L^2}^{3/2}\|\Delta u\|_{L^2}^{1/2}.
\]
Therefore
\[
  |\ip{R(u,\nabla u)}{-\Delta u}|
  \le CC_S\|\nabla u\|_{L^2}^{3/2}\|\Delta u\|_{L^2}^{3/2}
  \le\frac{\nu}{4}\|\Delta u\|_{L^2}^2+C_4E^3.
\]

\smallskip\noindent\textit{Combining.}
\[
  |\ip{(u\cdot\nabla)u}{-\Delta u}|
  \le\Bigl(\tfrac{\nu}{8}+\tfrac{\nu}{8}+\tfrac{\nu}{4}\Bigr)
     \|\Delta u\|_{L^2}^2+C_4E^3
  =\frac{\nu}{2}\|\Delta u\|_{L^2}^2+C_4E^3.
\]
Since $-\tfrac{\gamma}{4}\Tcal(u)\le 0$:
\[
  \frac{d}{dt}E+\frac{\nu}{2}\|\Delta u\|_{L^2}^2
  \le C_4E^3-\frac{\gamma}{4}\Tcal(u)\le C_4E^3.
  \qquad\square
\]
\end{proof}

\begin{remark}[Role of each case]
Case~I supplies the factored inequality used in the coercive bootstrap
(Proposition~\ref{prop:factored} and Theorem~\ref{thm:A}).
Case~II, where energy is too dispersed for $\Qsix$ to dominate, is
controlled by the Galerkin continuation argument in
Theorem~\ref{thm:gwp} (Step~3) via the $H^m$ estimate and the BKM
criterion; it does not require the factored form.
Global existence for the augmented system therefore holds
unconditionally.  The condition~[SND] is required only for the
de-augmentation step in Theorem~\ref{thm:A}.
\end{remark}

%% ─────────────────────────────────────────────────────────────────────────
\section{The Augmented System}
\label{sec:augmented}
%% ─────────────────────────────────────────────────────────────────────────

\subsection{Full formulation}
\begin{equation}\label{eq:aug}
  \partial_t u+(u\cdot\nabla)u+\nabla p
  =\nu\Delta u-\lH[u]\Delta u+\Qsix[u],
  \quad\nabla\cdot u=0,\quad u(\cdot,0)=u_0\in H^m,\;m\ge 3.
\end{equation}
Effective viscosity $\nu_{\rm eff}(t)=\nu+\lH[u(t)]\ge\nu>0$ uniformly.

\subsection{Augmented energy functional}
\begin{equation}\label{eq:EQ}
  E_Q[u]:=E[u]+\Phi[u],\quad\Phi[u]:=\tfrac{\alpha}{2}\|\nabla\psi[u]\|_{L^2}^2,
\end{equation}
\[
  \frac{d}{dt}E_Q\le-D[u],\quad
  D[u]:=\nu\|\nabla u\|_{L^2}^2+\lH\|\nabla u\|_{L^2}^2
  +\tfrac{\gamma}{4}\Tcal(u).
\]

%% ─────────────────────────────────────────────────────────────────────────
\section{Main Theorems for the Augmented System}
\label{sec:main}
%% ─────────────────────────────────────────────────────────────────────────

\begin{theorem}[Global well-posedness]\label{thm:gwp}
Let $u_0\in H^m(\R^3)$, $m\ge 3$, divergence-free, and let
$\gamma\ge\gamma_0:=8C_4E(0)$ where $C_4$ is the constant from
Lemma~\ref{lem:dichotomy}.  Then the augmented system~\eqref{eq:aug}
admits a unique solution
\[
  u\in C\bigl([0,\infty);H^m(\R^3)\bigr)
  \cap L^2_{\rm loc}\bigl([0,\infty);H^{m+1}(\R^3)\bigr),
\]
satisfying
\[
  \sup_{t\ge 0}\|u(t)\|_{H^m}^2
  +\int_0^\infty\|u(t)\|_{H^{m+1}}^2\,dt
  \le C\bigl(\|u_0\|_{H^m},\nu,\alpha,\gamma\bigr)<\infty.
\]
\end{theorem}

\begin{proof}
\textbf{Step~1: Local well-posedness via Carath\'eodory's theorem.}
The augmented system is quasilinear parabolic with $\nu_{\rm eff}\ge\nu>0$.
The operator $\Qsix$ is bounded $H^m\to H^{m-1}$ for a.e.\ $t$
(Lemma~\ref{lem:Q6}(v)).  The Galerkin finite-dimensional ODE has
right-hand side $f(t,u_n)$ that is (i)~locally Lipschitz in $u_n$ for
each fixed $t$, and (ii)~measurable and locally bounded in $t$ for fixed
$u_n$ (since $j_*(t)$ is measurable and piecewise constant).
These are the Carath\'eodory conditions (see~\cite{Hartman}, Ch.~II, Thm.~2.1),
which guarantee existence and uniqueness.
The $H^m$ estimate (Step~2) gives uniform bounds in $n$, and
Aubin--Lions gives a local strong solution
$u\in C([0,T^*];H^m)\cap L^2([0,T^*];H^{m+1})$.

\medskip\noindent
\textbf{Step~2: $H^m$ energy estimate.}
Apply $D^\alpha$ ($|\alpha|\le m$) to~\eqref{eq:aug} and test against
$D^\alpha u$.  By Lemma~\ref{lem:Q6}(iv), $D^\alpha$ and $m_{j_*}(D,t)$
commute, so
\[
  \ip{D^\alpha\Qsix[u]}{D^\alpha u}_{L^2}
  =-\gamma\int m_{j_*}(\xi,t)|\widehat{D^\alpha u}|^2\,d\xi\le 0.
\]
The $\Qsix$ term moves entirely to the left side as a non-negative
dissipative term; no $2^{4j_*}$ coefficient appears.
The Kato--Ponce commutator estimate~\cite{KP88} gives
$\|[D^\alpha,u\cdot\nabla]u\|_{L^2}\le C_m\|\nabla u\|_{L^\infty}\|u\|_{H^m}$.
With $\lH\ge 0$:
\[
  \frac{d}{dt}\|u\|_{H^m}^2+\nu\|u\|_{H^{m+1}}^2
  \le C\|\nabla u\|_{L^\infty}\|u\|_{H^m}^2.
\]
By Sobolev ($m\ge 3$): $\|\nabla u\|_{L^\infty}\le C\|u\|_{H^m}$, giving
$\frac{d}{dt}\|u\|_{H^m}^2\le C\|u\|_{H^m}^3$.
This ODE provides the local existence time and the continuation criterion.

\medskip\noindent
\textbf{Step~3: Global existence via bootstrap continuation.}
Bootstrap hypothesis~[BH]: suppose $E(t)\le M:=2E(0)$ on $[0,T]$.

\smallskip
\textit{Case~I holds on $[0,T]$:}
By Lemma~\ref{lem:dichotomy} Case~I,
$\frac{d}{dt}E\le\Tcal[C_4E-\gamma/4]$.
Under [BH] with $\gamma\ge\gamma_0=8C_4E(0)$:
\[
  C_4E(t)-\frac{\gamma}{4}
  \le C_4M-\frac{\gamma_0}{4}
  = 2C_4E(0)-2C_4E(0)=0.
\]
Since $\Tcal\ge 0$ and the bracket $\le 0$: $\frac{d}{dt}E\le 0$,
so $E(t)\le E(0)<M$.  Bootstrap closes strictly.

\smallskip
\textit{Case~II holds on $[0,T]$:}
By Lemma~\ref{lem:dichotomy} Case~II, $\frac{d}{dt}E\le C_4E^3$,
so the energy satisfies the ODE comparison $E(t)\le(E(0)^{-2}-2C_4t)^{-1/2}$,
which is finite on any interval $[0,T]$ with $T<1/(2C_4E(0)^2)$.
Combined with the uniform $H^m$ bound from Step~2 and the BKM criterion
($\|\nabla u\|_{L^\infty}\le C\|u\|_{H^m}<\infty$ implies
$\int_0^T\|\nabla u\|_{L^\infty}\,dt<\infty$), this precludes blow-up
and global existence follows.

\smallskip
In both cases the solution extends to all $t\ge 0$; higher Sobolev
norms follow by induction on $m$ using Step~2.
\end{proof}

\begin{theorem}[Uniform energy bound]\label{thm:energy}
Under Theorem~\ref{thm:gwp}: $E(t)\le E(0)$ for all $t\ge 0$ in
Case~I\textup{;} in general
\[
  \sup_{t\ge 0}E(t)\le C\bigl(\|u_0\|_{H^1},\nu,\gamma\bigr)<\infty.
\]
\end{theorem}
\begin{proof}
In Case~I, Step~3 gives $E(t)\le E(0)$ directly.
The uniform $H^m$ bound from Step~2 gives finite sup over all $T$;
the BKM argument of Step~3 extends to all time.
\end{proof}

\begin{theorem}[Vorticity control and BKM suppression]\label{thm:vort}
Under Theorem~\ref{thm:gwp}, the vorticity satisfies
\[
  \|\omega(t)\|_{L^\infty}\le M_\omega<\infty\quad\text{uniformly in }t\ge 0.
\]
\end{theorem}
\begin{proof}
$u\in C([0,\infty);H^m)$, $m\ge 3$.
Sobolev embedding $H^m\hookrightarrow C^{m-2}$ gives
$\|\nabla u\|_{L^\infty}\le C\|u\|_{H^m}\le C\sqrt{M}$ uniformly, so
$\|\omega\|_{L^\infty}\le C\|u\|_{H^m}=:M_\omega<\infty$.
Since BKM requires $\int_0^T\|\omega\|_{L^\infty}\,dt=\infty$ for
blow-up and $\|\omega\|_{L^\infty}$ is uniformly bounded, blow-up is
precluded.
\end{proof}

\begin{theorem}[Attractor convergence]\label{thm:attractor}
Under Theorem~\ref{thm:gwp}, there exists a compact set
$\Omega^*\subset H^1(\R^3)$ and constants $C_A,\delta>0$ such that
\[
  \operatorname{dist}_{H^1}(u(t),\Omega^*)\le C_A e^{-\delta t}
  \quad\text{for all }t\ge 0.
\]
\end{theorem}
\begin{proof}
The orbit $\{u(t)\}_{t\ge 0}$ is bounded in $H^m$ (Theorem~\ref{thm:energy});
by Rellich--Kondrachov its $H^1$-closure is compact.  Under~[SND] the
exponential rate $\delta=\frac{\gamma}{4}\frac{c_*}{2E(0)}$ follows from
$\frac{d}{dt}E\le-\delta E$.  Without~[SND], $H^1_{\rm weak}$-convergence
holds via the $H^m$ bound.
\end{proof}

%% ─────────────────────────────────────────────────────────────────────────
\section{Factored Energy Inequality Under [SND]}
\label{sec:prop2}
%% ─────────────────────────────────────────────────────────────────────────

\begin{proposition}[Factored active-shell energy inequality]
\label{prop:factored}
Assume~\textup{[SND]}: $\Tcal(u)\ge c_*\|\nabla u\|_{L^2}^2$ with
$c_*>0$ independent of time.  Then
\[
  \frac{d}{dt}E(t)\le\Tcal(u(t))\!\left[C_4 E(t)-\frac{\gamma}{4}\right].
\]
\end{proposition}
\begin{proof}
Under~[SND], Case~I of Lemma~\ref{lem:dichotomy} applies at every time
with $\sigma=c_*$.
\end{proof}

\begin{remark}
When $E(t)<\gamma/(4C_4)$ the bracket is negative and $\Tcal\ge 0$,
so $\frac{d}{dt}E\le 0$ regardless of $\Tcal$'s magnitude.
The choice $\gamma_0=8C_4E(0)$ ensures
$E(0)<\gamma_0/(4C_4)=2E(0)$, so the bracket is negative at $t=0$
and remains negative while $E(t)\le E(0)$, which is self-maintaining.
\end{remark}

%% ─────────────────────────────────────────────────────────────────────────
\section{The De-Augmentation Bridge and Theorem A}
\label{sec:deaug}
%% ─────────────────────────────────────────────────────────────────────────

\begin{definition}[Spectral Non-Dispersal Condition]\label{def:SND}
The family $\{u^\eps\}$ satisfies~\textup{[SND]} if
\[
  \inf_{\eps>0}\inf_{t\ge 0}
  \frac{\Tcal(u^\eps(t))}{\|\nabla u^\eps(t)\|_{L^2}^2}
  \ge c_*>0,
\]
with $c_*$ independent of $\eps$ and $t$.
\end{definition}

\noindent\textit{Physical meaning.}
[SND] asserts that the dominant three-shell band always captures a
positive fraction of gradient energy---the cascade never becomes
infinitely dispersed.  This is consistent with Kolmogorov K41
phenomenology but is not established by current PDE theory.

\begin{lemma}[De-augmentation convergence]\label{lem:deaug}
Under~\textup{[SND]}: $\Qsix(u^\eps)\to 0$ in $L^2(0,T;H^{-1}(\R^3))$
as $\eps\to 0$, and $u^\eps\to u$ strongly in $L^2(0,T;H^1_{\rm loc})$,
where $u$ solves the classical NS equations~\eqref{eq:NS}.
\end{lemma}
\begin{proof}
The symbol $m_{j_*}(\xi,t)=\gamma|\xi|^2\varphi(2^{-j_*}|\xi|)$.
For $\phi\in H^1$:
$|\ip{\Qsix(u^\eps)}{\phi}|\le C\gamma\cdot 2^{2j_*}\|u^\eps\|_{L^2}\|\phi\|_{H^1}$.
Under~[SND] and Theorem~\ref{thm:energy}, $2^{2j_*}$ is bounded
uniformly, so $\|\Qsix(u^\eps)\|_{H^{-1}}\le C\gamma\to 0$ as $\gamma\to 0$.
Aubin--Lions gives strong convergence; the limit satisfies classical NS
by passing to distributions.
\end{proof}

\begin{theorem}[Conditional global regularity of classical NS]
\label{thm:A}
Assume~\textup{[SND]}.  Then\textup{:}
\begin{enumerate}
  \item[\textup{(1)}]
    $\sup_{\eps>0}\sup_{t\ge 0}\|u^\eps(t)\|_{H^1}<\infty$.
  \item[\textup{(2)}]
    $u^\eps\to u$ strongly in $L^2_{\rm loc}(\R^3\times[0,\infty))$.
  \item[\textup{(3)}]
    $u\in C^\infty(\R^3\times(0,\infty))$ solves the classical NS
    equations.
\end{enumerate}
\end{theorem}
\begin{proof}
(1)~Under~[SND], Proposition~\ref{prop:factored} applies:
$\frac{d}{dt}E^\eps\le\Tcal[C_4E^\eps-\gamma/4]$.
With $\gamma=\gamma_0=8C_4E(0)$ fixed and
$E^\eps(0)\le E(0)<\gamma_0/(4C_4)=2E(0)$, the bracket at $t=0$ equals
$C_4E(0)-2C_4E(0)=-C_4E(0)<0$.
Since $\Tcal\ge 0$ and the bracket $<0$, $\frac{d}{dt}E^\eps\le 0$
initially.  The bracket remains negative whenever
$E^\eps\le E(0)<\gamma_0/(4C_4)$, which is maintained since $E^\eps$
is non-increasing.  Hence $E^\eps(t)\le E(0)$ for all $t\ge 0$,
uniformly in $\eps$.
(2)~Lemma~\ref{lem:deaug}.
(3)~$u\in L^\infty(0,\infty;H^1)$ solves classical NS; by Serrin's
regularity theorem~\cite{Ser62} and parabolic bootstrap,
$u\in H^k$ for all $k$, hence $u\in C^\infty$ for $t>0$.
\end{proof}

%% ─────────────────────────────────────────────────────────────────────────
\section{The Remaining Open Problem}
\label{sec:open}
%% ─────────────────────────────────────────────────────────────────────────

\begin{openproblem}[The Spectral Non-Dispersal Condition]
Prove or disprove: for any divergence-free $u_0\in H^1(\R^3)$, the
augmented solutions $\{u^\eps\}$ satisfy~\textup{[SND]}:
\[
  \inf_{\eps>0}\inf_{t\ge 0}
  \frac{\Tcal(u^\eps(t))}{\|\nabla u^\eps(t)\|_{L^2}^2}
  \ge c_*>0,
\]
with $c_*$ depending only on $\|u_0\|_{H^1}$ and $\nu$.
\end{openproblem}

\noindent\textbf{Non-circularity.}
[SND] does not assume regularity of the classical solutions.  It asks
about the augmented solutions, which are globally smooth by
Theorem~\ref{thm:gwp} for any fixed $\eps>0$.  A proof of~[SND]
combined with Theorem~\ref{thm:A} constitutes a complete proof of
global regularity for classical 3D Navier--Stokes.

\noindent\textbf{Possible approaches.}
(a)~Dynamical lower bounds on $\Tcal(u^\eps)$ from the augmented
equation's own structure: $\Qsix$ acts specifically on the dominant
shell, potentially reinforcing its dominance over time.
(b)~Profile decomposition: energy cannot disperse to infinite scales
under controlled dissipation.
(c)~Kolmogorov intermittency bounds establishing spectral concentration
in turbulent cascades.

%% ─────────────────────────────────────────────────────────────────────────
\section{Proof Status Summary}
\label{sec:status}
%% ─────────────────────────────────────────────────────────────────────────

\begin{center}
\begin{tabular}{lc}
\hline
\textbf{Result} & \textbf{Status}\\
\hline
$\Qsix$ definition (dominant shell $j_*$)            & Rigorous\\
Lemma~\ref{lem:Q6} --- coercive dissipation           & Rigorous\\
Lemma~\ref{lem:Q6}(iv) --- $H^m$ coercivity via Fourier commutativity & Rigorous\\
Lemma~\ref{lem:dichotomy} Preliminary --- $C_3$ explicit & Rigorous\\
Lemma~\ref{lem:dichotomy} Case~I --- corrected bootstrap chain & Rigorous\\
Lemma~\ref{lem:dichotomy} Case~II --- Bony paraproduct & Rigorous\\
Theorem~\ref{thm:gwp} --- augmented GWP, Carath\'eodory step 1 & Rigorous\\
Theorem~\ref{thm:gwp} --- $\gamma_0=8C_4E(0)$ algebra & Rigorous\\
Theorems~\ref{thm:energy}--\ref{thm:attractor}       & Rigorous\\
Proposition~\ref{prop:factored} --- factored inequality (under [SND]) & Rigorous\\
Theorem~\ref{thm:A} --- classical NS (conditional on [SND]) & Rigorous conditional\\
{[SND] Spectral Non-Dispersal Condition}              & \textbf{Open}\\
\hline
\end{tabular}
\end{center}

\noindent\textbf{One-sentence summary.}
The de-augmentation of the 3D Navier--Stokes regularity problem reduces
precisely to whether turbulent solutions satisfy~[SND]---a quantitative,
physically natural, mathematically precise statement about spectral
energy concentration that is currently open.

%% ─────────────────────────────────────────────────────────────────────────
\begin{thebibliography}{99}

\bibitem{BKM84}
J.~T.~Beale, T.~Kato, and A.~Majda,
\textit{Remarks on the breakdown of smooth solutions for the 3-D Euler
equations},
Comm.\ Math.\ Phys.\ \textbf{94} (1984), 61--66.

\bibitem{Bon81}
J.-M.~Bony,
\textit{Calcul symbolique et propagation des singularit\'es pour les
\'equations aux d\'eriv\'ees partielles non lin\'eaires},
Ann.\ Sci.\ \'Ecole Norm.\ Sup.\ \textbf{14} (1981), 209--246.

\bibitem{CF88}
P.~Constantin and C.~Foias,
\textit{Navier--Stokes Equations},
University of Chicago Press, 1988.

\bibitem{CKN82}
L.~Caffarelli, R.~Kohn, and L.~Nirenberg,
\textit{Partial regularity of suitable weak solutions of the
Navier--Stokes equations},
Comm.\ Pure Appl.\ Math.\ \textbf{35} (1982), 771--831.

\bibitem{CM78}
R.~Coifman and Y.~Meyer,
\textit{Au del\`a des op\'erateurs pseudo-diff\'erentiels},
Ast\'erisque \textbf{57} (1978).

\bibitem{ESS03}
L.~Escauriaza, G.~Seregin, and V.~\v{S}ver\'ak,
\textit{$L_{3,\infty}$-solutions of Navier--Stokes equations and
backward uniqueness},
Uspekhi Mat.\ Nauk \textbf{58} (2003), 3--44.

\bibitem{Hartman}
P.~Hartman,
\textit{Ordinary Differential Equations},
2nd ed., SIAM Classics in Applied Mathematics, 2002.

\bibitem{Hop51}
E.~Hopf,
\textit{\"Uber die Anfangswertaufgabe f\"ur die hydrodynamischen
Grundgleichungen},
Math.\ Nachr.\ \textbf{4} (1951), 213--231.

\bibitem{KP88}
T.~Kato and G.~Ponce,
\textit{Commutator estimates and the Euler and Navier--Stokes equations},
Comm.\ Pure Appl.\ Math.\ \textbf{41} (1988), 891--907.

\bibitem{Ler34}
J.~Leray,
\textit{Sur le mouvement d'un liquide visqueux emplissant l'espace},
Acta Math.\ \textbf{63} (1934), 193--248.

\bibitem{Ser62}
J.~Serrin,
\textit{On the interior regularity of weak solutions of the
Navier--Stokes equations},
Arch.\ Ration.\ Mech.\ Anal.\ \textbf{9} (1962), 187--195.

\bibitem{Tao16}
T.~Tao,
\textit{Finite time blowup for an averaged three-dimensional
Navier--Stokes equation},
J.\ Amer.\ Math.\ Soc.\ \textbf{29} (2016), 601--674.

\bibitem{Tem79}
R.~Temam,
\textit{Navier--Stokes Equations},
North-Holland, 1979.

\end{thebibliography}

\end{document}
